\section{From the certificates to the game}\label{sec:conclusion}
We now combine the three obligations identified in the introduction.
The next implication is mathematical; its finite premise is exactly the
input and certificate assertion of Proposition~\ref{prop:finite}.

\begin{theorem}[Certificate-to-game implication]\label{thm:bridge}
If the finite exclusion and input checks in
Proposition~\ref{prop:finite} hold, then no five-color winning strategy
exists on $C_5$.
\end{theorem}
\begin{proof}
Suppose that $f$ is a five-color winning strategy. By
Proposition~\ref{prop:coverage}, a normalized winning strategy $g$ obtained
from it has count vector $a\in\calR_n$, with $17\leq n\leq25$, and
satisfies $g_1(5,5)=1$. There are two cases.

If this parent is a direct terminal branch,
Proposition~\ref{prop:input} extends the strategy's base assignment to a
satisfying assignment of its actual formula $F_a$. The checked refutation
of $F_a$ makes this impossible by Lemma~\ref{lem:rup}.

If the parent is refined, Lemma~\ref{lem:patterns} first applies a
common permutation of the four non-special colors to $g$, placing
its diagonal answers in some $t\in\calT$. This permutation preserves the
parent count vector and the condition that the first diagonal answer is 1.
Apply the input-preservation result
\emph{to this transformed strategy}: choose new auxiliary assignments
for its five counter components. There is no need to transport a previous
auxiliary assignment through the color permutation or to assume that the
counter's syntax is color-equivariant. The resulting satisfying assignment
of $F_{a,t}$ contradicts its checked refutation.

Every parent is either direct or refined, and every refined parent has
all 51 required subcases. Thus the assumed strategy has been excluded.
\end{proof}

\begin{proof}[Proof of Theorem~\ref{thm:main}]
The complete input audit and certificate replay of
Section~\ref{sec:certificates} verify the finite premise of
Theorem~\ref{thm:bridge}, excluding five colors. Proposition~\ref{prop:four}
gives a winning four-color strategy. A winning strategy with more than
five colors would restrict to one with five by Lemma~\ref{lem:palette}.
Consequently $\HG(C_5)=4$.
\end{proof}

\subsection{Why weight 25 is the final obstruction}
The exclusions below 25 have a useful interpretation before the last
weight is excluded. This observation is a consequence of those exclusions,
not an assumption in their Boolean inputs.

\begin{corollary}[The last remaining count]\label{cor:balance}
Assume the exclusions at weights 17--24 hold. Any five-color winning
strategy would satisfy
\[
S(1)=S(2)=S(3)=S(4)=S(5)=25
\]
and would have no entry guessing a neighbor's color. The exclusion at
weight 25 then contradicts this necessary condition.
\end{corollary}
\begin{proof}
Fix a color $z$ in the original strategy. Use the repair of
Lemma~\ref{lem:normalization} that preserves the indicators for this
particular $z$. The bound of Proposition~\ref{prop:seventeen} and the
exclusions at weights 17--24 imply $S(z)\geq25$: otherwise this repaired
strategy could be transformed into one of the excluded branches.
Repeat this argument for each color. The repairs may differ with $z$;
each inequality concerns the corresponding count in the \emph{same
original strategy}, because that count was preserved.

Now apply Lemma~\ref{lem:total} to the original strategy:
$\sum_z S(z)\leq125$. All five lower bounds must be equalities, and every
table entry must contribute to the sum. An entry guessing a neighbor's
color contributes zero and is therefore impossible. The weight-25
exclusion eliminates the remaining possibility.
\end{proof}

The equality is a total over five players for each fixed color. It does
not assert that each player guesses each color five times. None of the
input formulas used this balance conclusion or a lower bound already
obtained from another weight, so the reasoning has no circular dependence.

\subsection{Computational evidence and access}
The sealed computational deposit has version
\begin{quote}\small\raggedright
\path{c5-reviewed-evidence-20261027-v2}.
\end{quote}
It contains selected source code,
English reproduction instructions, the two-stage review archive and all
nine complete evidence archives. The precise inventories and commands
are given in Appendix~\ref{app:reproduction}; source fingerprints and
actual run records are identified in Appendix~\ref{app:audit}.
Saved-evidence verification does not require repeating the original SAT
search. It requires checking branch coverage and input preservation, and
then replaying all terminal refutations.

The selected files are staged in a private repository. They will be made
public when this preprint is announced on arXiv, at the fixed release address
\begin{quote}\small\raggedright
\url{https://github.com/Lshiqi67/proper-hat-guessing-c5-reproducibility/releases/tag/c5-arxiv-v1-evidence-20261027}.
\end{quote}
Until that coordinated release, access is restricted. No archive DOI has
been assigned. Hashes identify versions and detect byte changes; they do
not replace access to the files or validation of their mathematical meaning.

The sealed deposit includes the 26 October audited manuscript that was
examined in the two-stage review. This submission candidate derives from
the 27 October exposition revision and is supplied separately, with its
source package and version correspondence. The definitions, branch inputs
and certificates are unchanged; preparing this candidate neither alters
the sealed evidence nor constitutes another full replay.

\subsection{Assistance and further mathematical work}
OpenAI Codex assisted with mathematical exploration, program development,
local computations, evidence organization and manuscript preparation.
The solver, input auditor and proof checker serve separate purposes,
but their development occurred within the same AI-assisted workflow.
External mathematical scrutiny and a separately reviewed or formally
verified checker remain useful further checks of the argument.

A mathematical question left by this proof is whether the finite
exclusion can be replaced by a shorter explanation of the incompatibility
between the zero-, one- and two-occurrence families. The witness in
Appendix~\ref{app:witness} shows that the two-occurrence restriction alone
permits total count 17. A shorter proof must therefore use additional
compatibility information, rather than simply strengthen that local
bound for the same restricted family.
